% Slides boundary extraction: N in {5,7,8} reinterpretation
% Primary sources:
%   Slides plain text: /tmp/automorphic-slides.txt (953 lines)
%   Paper PDF: /Users/raeez/Downloads/raeez.lorgat.automorphic-corrections.pdf
% Cross-citations: Mukai1988, GritsenkoNikulin1998, GritsenkoClery2008,
%   ChengDuncanHarvey2014, Lorgat2020.

\providecommand{\kBKM}{\kappa_{\mathrm{BKM}}}
\providecommand{\kch}{\kappa_{\mathrm{ch}}}
\providecommand{\kcat}{\kappa_{\mathrm{cat}}}
\providecommand{\kfib}{\kappa_{\mathrm{fiber}}}
\providecommand{\gDelta}{\mathfrak{g}_{\Delta_5}}
\providecommand{\ClaimStatusConjectured}{[Conjectured]}
\providecommand{\ClaimStatusTheorem}{[Theorem]}
\providecommand{\ClaimStatusAxiomatic}{[Axiomatic]}

\section*{The boundary $N \in \{5,7,8\}$ examined against the slides}

\paragraph{The eight Gritsenko--Cl\'ery dd-modular forms (slide line 47--54).}
\[
\begin{array}{l@{\;}l@{\;}l}
  M_5(\Gamma_1,\nu_2)=\Delta_5 & M_2(\Gamma_2,\nu_4) & M_3(\Gamma_1(2),\nu_2)\\
  M_1(\Gamma_3,\nu_6) & M_2(\Gamma_1(3),\nu_2) & \\
  M_{1/2}(\Gamma_4,\nu_8) & M_{3/2}(\Gamma_1(4),\nu_4) & M_1(\Gamma_2(2),\nu_4)
\end{array}
\]
\ClaimStatusTheorem, \cite{GritsenkoClery2008}.

\paragraph{Factual split by weight.} Six integer-weight forms
($\Delta_5$, $M_2(\Gamma_2)$, $M_3(\Gamma_1(2))$, $M_1(\Gamma_3)$,
$M_2(\Gamma_1(3))$, $M_1(\Gamma_2(2))$) and two half-integer-weight
forms ($M_{1/2}(\Gamma_4,\nu_8)$, $M_{3/2}(\Gamma_1(4),\nu_4)$). The
two half-integer forms sit at paramodular level $t=4$, \emph{not} at
$N \in \{5,7,8\}$. A 5+3 split keyed to $\{5,7,8\}$ is inconsistent
with slide line~54.

\paragraph{The slides' own case-split (line 943--953).}
\[
N \in \{1,2,3,5,7\}\;:\;1\le s,r\le N{-}1,\;0\le k\le N{-}1;\quad
N=4,6,8\;:\;\text{separate sectors}.
\]
$N=8$ appears twice (lines 951 and 953) — two slide-pages, two
orbit-sectors at the Mukai bound.

\paragraph{Boundary interpretation dictionary.}
\begin{tabular}{@{}l@{\;}l@{\;}l@{}}
\hline
$N$ & Slide treatment & Programme reinterpretation \\ \hline
$5$ & Inside $\{1,2,3,5,7\}$ uniform block (line 943) &
  \emph{Not boundary.} Sits with $\{1,2,3,7\}$; order-5 symplectic
  $g_5 \in M_{23}$ acts with $1+4$ fixed-lattice type on $H^*(K3)$;
  Mukai class $5A$.\\
$7$ & Inside $\{1,2,3,5,7\}$ uniform block (line 943) &
  \emph{Not boundary.} Sits with $\{1,2,3,5\}$; order-7 symplectic
  Mukai class $7A$; twisted K3 elliptic genus is holomorphic Jacobi
  of weight $0$, index $1$ at level $7$ by
  \cite{ChengDuncanHarvey2014} $M_{24}$ moonshine tables.\\
$8$ & Separate, doubled (lines 951, 953) &
  \emph{Is boundary.} Saturates Mukai's bound $N \le 8$
  \cite{Mukai1988} Theorem~0.3; two orbit-sectors = $M_{24}$
  classes $8A, 8B$ both restricting to order-$8$ symplectic on a
  Mukai-lattice embedding.\\
\hline
\end{tabular}

\paragraph{The Wave~D\,2/D\,3 5+3 split, reinterpreted.} The
structural 5+3 is \emph{not} at $N \in \{5,7,8\}$ but at
\emph{$N \in \{4,6,8\}$ separate} vs. \emph{$N \in \{1,2,3,5,7\}$
uniform}. The discriminant is the $2$-divisibility of $N$: when
$\mathbb{Z}/N\mathbb{Z}$ has a non-trivial order-$2$ subgroup
(i.e.\ $N$ even), the $(S \times E)/\mathbb{Z}/N$ quotient has fixed
fibres over the $2$-torsion of $E$, and the orbit-index set degenerates
($r,k$ become absorbed into $s$). For odd $N$ or $N$ prime to $2$, the
free-action index set $(s,r,k)$ with $s,r \in \{1,\dots,N{-}1\}$ is
generic. The $\{4,6,8\}$ separation is not ``boundary / non-compact
CHL'' but \emph{compositeness-with-$2$} of the orbifold group.

\paragraph{Half-integer weight at $N=4$, not $N=8$.}
\ClaimStatusTheorem. The two half-integer-weight dd-modular forms
$M_{1/2}(\Gamma_4,\nu_8)$ and $M_{3/2}(\Gamma_1(4),\nu_4)$ occur at
paramodular level $t=4$. The paramodular $\Gamma_4$ admits
$\nu_8$-multiplier characters whose metaplectic double-cover arguments
lift $\eta$-type Jacobi theta-series to half-integer-weight Siegel
forms; \cite{GritsenkoClery2008} Table~1 lists these explicitly. No
half-integer weight appears at $N \in \{5,7,8\}$: the $N=8$ GC form is
not in the slides' eight-form list (slide line~54); the $N \in \{5,7\}$
cases contribute via $\ZLh$ products of integer-weight forms only.

\paragraph{Mukai bound at $N=8$.}
\ClaimStatusTheorem, \cite{Mukai1988} Theorem~0.3. Every finite group
$G$ of symplectic automorphisms of a K3 surface embeds in $M_{23}$; the
maximal cyclic order is $N=8$. The slide's doubled ``Case $N=8$''
(lines 951, 953) reflects the existence of two $M_{24}$-classes
$8A,8B$ in Mukai's classification with distinct Niemeier embeddings.
The Mukai lattice is standard (not non-standard); no non-Niemeier
umbral orbit appears at $N \le 8$ per \cite{ChengDuncanHarvey2014}.

\paragraph{Resolution of the hidden-observation test.}
\ClaimStatusTheorem. The hypothesis ``$N \in \{5,7,8\}$ corresponds
to genuinely different Mukai-group-theoretic structures (non-Niemeier
umbral, non-standard Mukai lattice)'' is \emph{falsified}. All
$N \le 8$ symplectic-automorphism orbits fit Mukai's Niemeier-embedded
classification. The true boundary phenomenon at $N=8$ is
\emph{Mukai-bound saturation}, not non-Niemeier exoticism. $N=5$ and
$N=7$ are not boundary cases; the slides place them in the uniform
block.

\paragraph{Invariants at $N=8$.}
$\kch(K3 \times E) = 0$ (K\"unneth, compact CY$_3$). At the
$(S \times E)/\mathbb{Z}/8$ quotient $X$:
$\kcat(X) = \chi(\mathcal{O}_X) = 0$ (CY$_3$, total-space K\"unneth
survives quotients preserving CY form). $\kBKM$ for the putative
$N=8$ Borcherds correction: if $\Phi_8$ denotes the Gritsenko
$\Phi_N$-series \cite{GritsenkoNikulin1998}, $\kBKM(\Phi_8) =
c_8(0)/2$; the universal identity
$\kBKM(\Phi_N) = c_N(0)/2$ (Vol~III main theorem) extends to $N=8$
only if the twisted K3 elliptic genus at Mukai class $8A$ or $8B$ is
a Jacobi form of weight $0$ index $1$; this is
\ClaimStatusConjectured\ per \cite{Lorgat2020} Conjecture~1 step~S5.

\paragraph{Summary.} The boundary $N \in \{5,7,8\}$ posed in
Wave~D\,2/D\,3 is not the structural boundary exhibited by the
slides. The slides' split is $\{1,2,3,5,7\}$ uniform vs.\
$\{4,6,8\}$ separate-sector; the half-integer-weight forms sit at
$N=4$ (paramodular level), not at $\{5,7,8\}$; the Mukai bound
saturates only at $N=8$, and there via two $M_{24}$-classes
$8A,8B$, not via non-Niemeier exoticism. Programme reinterpretation:
replace ``$N \in \{5,7,8\}$ boundary'' by ``$N \in \{4,6,8\}$
even-compositeness separation plus Mukai saturation at $N=8$''.
